\chapter{Grafy losowe i metoda probabilistyczna}
\label{ch:tg-random-graphs}

Twierdzenia Mantela, Turána i Ramseya z rozdziału
\ref{ch:tg-extremal-ramsey} dotyczyły wszystkich grafów
spełniających określone założenia. Teraz będziemy
losować grafy, aby obliczać prawdopodobieństwa ich
własności i wykazywać istnienie grafów
o określonej strukturze. Porównamy modele $G(n,p)$ i $G(n,m)$, przeanalizujemy
liczby krawędzi, trójkątów i wierzchołków izolowanych,
a następnie udowodnimy progi pojawiania się ustalonych
podgrafów i spójności. Na koniec zastosujemy metodę
probabilistyczną do liczb Ramseya.

\section{Model Erdősa--Rényiego i zmienne losowe}
\label{sec:tg-random-model}

\begin{definicja}[Losowy graf $G(n,p)$]
\label{def:tg-random-gnp}
Niech $n\geq0$ będzie liczbą całkowitą,
$p\in[0,1]$ oraz
$V=\{1,\ldots,n\}$. Model $G(n,p)$
jest grafem prostym nieskierowanym
na zbiorze $V$, w którym każda
z $\binom n2$ możliwych krawędzi
jest wybierana niezależnie
z prawdopodobieństwem $p$.

Dokładniej, dla każdej pary
$e\in\binom V2$ rozważamy niezależną
zmienną Bernoulliego $I_e$
o wartości $1$ z prawdopodobieństwem
$p$ i $0$ z prawdopodobieństwem
$1-p$. Zbiór krawędzi to
$E=\{e:I_e=1\}$.
\end{definicja}

Dla konkretnego grafu $H$ na
oznaczonym zbiorze $V$ zachodzi
\[
\mathbb P(G(n,p)=H)
=p^{|E(H)|}(1-p)^{\binom n2-|E(H)|},
\]
z naturalną konwencją dla
$p=0$ i $p=1$. Model ten
nie nakłada warunku spójności
ani stałej liczby krawędzi.
Jeżeli parametr zależy od $n$,
zapisujemy $p=p_n$.

\begin{definicja}[Zdarzenie zachodzące z prawdopodobieństwem
dążącym do jedności]
\label{def:tg-random-whp}
Niech $A_n$ będzie zdarzeniem
dotyczącym grafu $G(n,p_n)$.
Mówimy, że $A_n$ zachodzi
\emph{z prawdopodobieństwem
dążącym do jedności}, jeśli
$\lim_{n\to\infty}\mathbb P(A_n)=1$.
Analogicznie własność zanika
z prawdopodobieństwem dążącym
do jedności, jeśli jej
prawdopodobieństwo dąży do zera.
\end{definicja}

\begin{stwierdzenie}[Krawędzie i stopnie]
\label{prop:tg-random-edges-degrees}
Niech $n\geq1$ będzie liczbą całkowitą i $p\in[0,1]$.
Dla grafu $G(n,p)$
liczba krawędzi $M=|E|$
ma rozkład dwumianowy
$\operatorname{Bin}(\binom n2,p)$.
W szczególności
\[
\mathbb E M=\binom n2p,
\qquad
\operatorname{Var}(M)=\binom n2p(1-p).
\]
Dla każdego $v\in V$
stopień $d_G(v)$ ma rozkład
$\operatorname{Bin}(n-1,p)$,
a więc jego wartość oczekiwana
wynosi $(n-1)p$, a wariancja
$(n-1)p(1-p)$.
\end{stwierdzenie}
\begin{proof}
Mamy $M=\sum_{e\in\binom V2}I_e$,
a $d_G(v)=\sum_{e\ni v}I_e$.
Pierwsza suma ma $\binom n2$
składników, druga $n-1$,
a wszystkie składniki
w każdej sumie są niezależnymi
zmiennymi Bernoulliego
z parametrem $p$.
Wzory wynikają z liniowości
wartości oczekiwanej
i addytywności wariancji
dla zmiennych niezależnych.
\end{proof}

\begin{wniosek}[Koncentracja liczby krawędzi]
\label{cor:tg-random-edge-concentration}
Jeżeli $p_n\in[0,1]$ oraz
$\mu_n=\binom n2p_n\to\infty$,
to dla każdego $\varepsilon>0$
\[
\mathbb P\bigl(
|M-\mu_n|\geq\varepsilon\mu_n
\bigr)\longrightarrow0.
\]
\end{wniosek}
\begin{proof}
Dla każdego $n$ z $\mu_n>0$
nierówność Czebyszewa daje
\[
\mathbb P(|M-\mu_n|\geq\varepsilon\mu_n)
\leq\frac{\operatorname{Var}(M)}
{\varepsilon^2\mu_n^2}
=\frac{1-p_n}{\varepsilon^2\mu_n}
\leq\frac1{\varepsilon^2\mu_n}
\longrightarrow0.
\]
\end{proof}

\begin{przyklad}[Mały model probabilistyczny]
\label{ex:tg-random-four}
Dla $G(4,\tfrac12)$ istnieje
sześć możliwych krawędzi,
z których każda jest wybierana
z prawdopodobieństwem $\tfrac12$.
Wszystkie $2^6=64$ grafy
na oznaczonym zbiorze
czterech wierzchołków
są jednakowo prawdopodobne.
Otrzymujemy $\mathbb EM=3$
oraz $\operatorname{Var}(M)=\tfrac32$.
Prawdopodobieństwo braku
jakiejkolwiek krawędzi
wynosi $2^{-6}$.
\end{przyklad}

\begin{figure}[htbp]
\centering
\begin{tikzpicture}[>=Stealth,scale=.92]
\foreach \i in {1,...,6}{
  \pgfmathsetmacro{\ang}{90-60*(\i-1)}
  \node[tgVertex] (v\i) at (\ang:1.85) {$\i$};
}
\draw[tgTree] (v1)--(v2)--(v3);
\draw[tgTree] (v3)--(v4);
\draw[tgTree] (v4)--(v6);
\draw[tgTree] (v6)--(v1);
\draw[tgTree] (v2)--(v5);
\draw[tgTree] (v1)--(v3);
\end{tikzpicture}
\caption{Jedna z możliwych realizacji
modelu $G(6,p)$. Rysunek przedstawia
konkretny graf, a nie średnią
ani strukturę typową dla każdego $p$.}
\label{fig:tg-random-realization}
\end{figure}

\section{Model o ustalonej liczbie krawędzi i sprzężenie modeli}
\label{sec:tg-random-gnm}

\begin{definicja}[Model $G(n,m)$]
\label{def:tg-random-gnm}
Niech $n,m$ będą liczbami całkowitymi, $n\geq0$,
$0\leq m\leq N=\binom n2$. Model $G(n,m)$
jest jednostajnym rozkładem na zbiorze grafów prostych
na $\{1,\ldots,n\}$ mających dokładnie $m$ krawędzi.
Każdy z $\binom Nm$ takich grafów ma prawdopodobieństwo
$1/\binom Nm$.
\end{definicja}

\begin{stwierdzenie}[Związek modeli]
\label{prop:tg-random-conditional-gnm}
Jeżeli $M$ jest liczbą krawędzi w $G(n,p)$ oraz
$\mathbb P(M=m)>0$, to rozkład $G(n,p)$ warunkowo
względem zdarzenia $\{M=m\}$ jest równy rozkładowi
$G(n,m)$. Dla dowolnego ustalonego zbioru $F$
składającego się z $h$ różnych możliwych krawędzi,
$0\leq h\leq N$, zachodzi
\[
\mathbb P(F\subseteq E(G(n,m)))=
\begin{cases}
\binom{N-h}{m-h}/\binom Nm,&h\leq m,\\
0,&h>m.
\end{cases}
\]
\end{stwierdzenie}
\begin{proof}
Każdy graf o $m$ krawędziach ma w modelu $G(n,p)$
to samo prawdopodobieństwo $p^m(1-p)^{N-m}$.
Po podzieleniu przez $\mathbb P(M=m)=
\binom Nm p^m(1-p)^{N-m}$ otrzymujemy rozkład
jednostajny. Przy $p=0,1$ jedyne zdarzenie
warunkujące o dodatnim prawdopodobieństwie
odpowiada jedynemu możliwemu grafowi.
Jeśli $h\leq m$, graf zawierający $F$ powstaje
przez wybór pozostałych $m-h$ krawędzi spośród
$N-h$ krawędzi poza $F$. Jeśli $h>m$, taki graf
nie istnieje.
\end{proof}

\begin{przyklad}[Zależność krawędzi w modelu $G(n,m)$]
Dla $N\geq2$ oraz dwóch różnych możliwych krawędzi
$e,f$ mamy
$\mathbb P(e\in E)=m/N$ oraz
$\mathbb P(e,f\in E)=m(m-1)/(N(N-1))$.
Stąd kowariancja ich funkcji charakterystycznych
wynosi $-m(N-m)/(N^2(N-1))$. Dla $0<m<N$ jest
ujemna, a więc wybory krawędzi nie są niezależne,
w odróżnieniu od modelu $G(n,p)$.
\end{przyklad}

\begin{definicja}[Własność rosnąca]
\label{def:tg-random-increasing}
Własność grafów na ustalonym zbiorze wierzchołków
nazywamy rosnącą, jeśli z jej zachodzenia dla grafu
$(V,E)$ wynika jej zachodzenie dla każdego grafu
$(V,E')$ z $E\subseteq E'$. Przykładami są spójność
i zawieranie ustalonego podgrafu.
\end{definicja}

\begin{stwierdzenie}[Sprzężenie monotoniczne]
\label{prop:tg-random-coupling}
Dla $0\leq p\leq q\leq1$ i ustalonego $n$
można zdefiniować na jednej przestrzeni
probabilistycznej grafy $G_p,G_q$ o rozkładach
$G(n,p),G(n,q)$ tak, że $E(G_p)\subseteq E(G_q)$
z prawdopodobieństwem jeden. W szczególności
dla własności rosnącej $\mathcal A$ zachodzi
$\mathbb P(G_p\in\mathcal A)\leq\mathbb P(G_q\in\mathcal A)$.
\end{stwierdzenie}
\begin{proof}
Każdej możliwej krawędzi $e$ przypiszmy niezależną
zmienną $U_e$ o rozkładzie jednostajnym na $[0,1]$.
Dla $0\leq x\leq1$ włączamy $e$ do $G_x$ wtedy,
gdy $U_e<x$. Zdarzenia te są niezależne
dla różnych $e$ i mają prawdopodobieństwo $x$.
Ponadto $U_e<p$ implikuje $U_e<q$ dla $p\leq q$.
Zawieranie zbiorów krawędzi daje punktową
nierówność funkcji charakterystycznych
własności rosnącej, a po wzięciu wartości
oczekiwanych otrzymujemy tezę.
\end{proof}

\section{Metoda pierwszego i drugiego momentu}
\label{sec:tg-random-moments}

\begin{lemat}[Pierwszy moment]
\label{lem:tg-random-first-moment}
Niech $X$ będzie nieujemną
zmienną losową o skończonej
wartości oczekiwanej.
Wtedy
$\mathbb P(X\geq1)\leq\mathbb EX$.
Jeżeli dodatkowo $X$ jest
całkowitoliczbowa i
$\mathbb EX<1$, to
$\mathbb P(X=0)>0$.
\end{lemat}
\begin{proof}
Nierówność Markowa
$\mathbb P(X\geq1)\leq\mathbb EX$
wynika bezpośrednio
z $\mathbb EX\geq
\mathbb E(X\chi_{\{X\geq1\}})
\geq\mathbb P(X\geq1)$.
Jeśli $X$ jest nieujemna
i całkowitoliczbowa,
zdarzenia $X=0$ oraz
$X\geq1$ są dopełniające,
więc przy $\mathbb EX<1$
dostajemy
$\mathbb P(X=0)\geq1-\mathbb EX>0$.
\end{proof}

\begin{lemat}[Drugi moment i nieistnienie zera]
\label{lem:tg-random-second-moment}
Niech $X$ będzie nieujemną
zmienną losową o skończonej
wariancji, dla której
$\mu=\mathbb EX>0$.
Wtedy
\[
\mathbb P(X=0)
\leq\frac{\operatorname{Var}(X)}{\mu^2}.
\]
W szczególności, jeżeli
dla ciągu $X_n$ zachodzi
$\mathbb EX_n>0$ oraz
$\operatorname{Var}(X_n)/
(\mathbb EX_n)^2\to0$,
to $\mathbb P(X_n>0)\to1$.
\end{lemat}
\begin{proof}
Zdarzenie $\{X=0\}$ zawiera
się w $\{|X-\mu|\geq\mu\}$.
Stosujemy nierówność Czebyszewa
i przechodzimy do granicy.
\end{proof}

Zaletą metody pierwszego
momentu jest to, że
wystarczy obliczyć
średnią liczbę
zabronionych konfiguracji.
Metoda drugiego momentu
uwzględnia natomiast
zależności między
wystąpieniami tych
konfiguracji.

\section{Trójkąty i próg ich pojawiania się}
\label{sec:tg-random-triangles}

\begin{stwierdzenie}[Średnia i wariancja liczby trójkątów]
\label{prop:tg-random-triangle-moments}
Niech $n\geq3$, $p\in[0,1]$,
a $X$ oznacza liczbę
nieuporządkowanych trójek
wierzchołków tworzących
trójkąt w $G(n,p)$.
Wówczas
\[
\mathbb EX=\binom n3p^3
\]
oraz
\[
\operatorname{Var}(X)
=\binom n3p^3(1-p^3)
+2\binom n2\binom{n-2}2p^5(1-p).
\]
\end{stwierdzenie}
\begin{proof}
Niech $\mathcal T=\binom V3$
i niech $J_T$ będzie funkcją charakterystyczną
zdarzenia obecności trójkąta
na trójce $T$.
Wtedy $X=\sum_{T\in\mathcal T}J_T$,
a $\mathbb EJ_T=p^3$,
ponieważ konieczna jest
obecność wszystkich
trzech niezależnie
wybieranych krawędzi.

Dla różnych trójek $T,U$
funkcje charakterystyczne są niezależne,
jeśli trójkąty
nie mają wspólnej krawędzi;
mogą przy tym dzielić
jeden wierzchołek.
Jeśli mają wspólną
krawędź, ich jednoczesne wystąpienie
wymaga obecności pięciu
różnych krawędzi, więc
$\operatorname{Cov}(J_T,J_U)
=p^5-p^6=p^5(1-p)$.
Liczba \emph{nieuporządkowanych}
par różnych trójkątów
mających wspólną
krawędź to
$\binom n2\binom{n-2}2$:
najpierw wybieramy
wspólną krawędź,
potem dwa różne
wierzchołki uzupełniające.
Korzystając z tożsamości
$\operatorname{Var}(\sum J_T)
=\sum\operatorname{Var}(J_T)
+2\sum_{T<U}\operatorname{Cov}(J_T,J_U)$,
otrzymujemy wzór.
\end{proof}

\begin{twierdzenie}[Próg występowania trójkąta]
\label{thm:tg-random-triangle-threshold}
Niech $p_n\in[0,1]$.
Jeżeli $np_n\to0$, to
$\mathbb P(G(n,p_n)\text{ zawiera }K_3)\to0$.
Jeżeli $np_n\to\infty$, to
$\mathbb P(G(n,p_n)\text{ zawiera }K_3)\to1$.
Zatem funkcja $n^{-1}$
wyznacza próg pojawiania się
trójkąta w powyższym sensie.
\end{twierdzenie}
\begin{proof}
Jeśli $np_n\to0$, to
$\mathbb EX=\binom n3p_n^3
\leq(np_n)^3/6\to0$.
Pierwszy moment
\ref{lem:tg-random-first-moment}
daje $\mathbb P(X>0)\to0$.

Jeśli $np_n\to\infty$,
to dla dużych $n$
wartość $\mu_n=\binom n3p_n^3$
jest dodatnia i rośnie
do nieskończoności.
Ze wzoru na wariancję
i $\binom n3\asymp n^3$
dostajemy
\[
\frac{\operatorname{Var}(X)}{\mu_n^2}
\leq\frac1{\mu_n}
+\frac{2\binom n2\binom{n-2}2}
{\binom n3^2\,p_n}
=O\left(\frac1{(np_n)^3}
+\frac1{n^2p_n}\right)\to0.
\]
Ostatnia granica wynika
z $np_n\to\infty$.
Lemat
\ref{lem:tg-random-second-moment}
daje $\mathbb P(X=0)\to0$,
a więc $\mathbb P(X>0)\to1$.
\end{proof}

\begin{uwaga}[Co oznacza próg]
Twierdzenie nie stwierdza,
że dla pojedynczego $n$
graf mający $p>1/n$
musi zawierać trójkąt.
Granice dotyczą ciągów
prawdopodobieństw
i zachowania
przy $n\to\infty$.
W obszarze $p_n\asymp1/n$
prawdopodobieństwo
może pozostać
oddalone zarówno od zera,
jak i od jedności.
\end{uwaga}
\section{Wierzchołki izolowane i próg spójności}
\label{sec:tg-random-connectivity}

\begin{stwierdzenie}[Liczba wierzchołków izolowanych]
\label{prop:tg-random-isolates}
Niech $n\geq2$, $p\in[0,1]$
i niech $Y$ będzie liczbą
wierzchołków stopnia zero
w $G(n,p)$. Wówczas
\[
\mathbb EY=n(1-p)^{n-1}
\]
oraz
\[
\operatorname{Var}(Y)=
nq(1-q)+n(n-1)p(1-p)^{2n-3},
\qquad q=(1-p)^{n-1}.
\]
\end{stwierdzenie}
\begin{proof}
Dla $v\in V$ niech $J_v$
będzie funkcją charakterystyczną zdarzenia, że
$v$ jest izolowany. Zdarzenie $J_v=1$
wymaga braku wszystkich
$n-1$ krawędzi przy $v$,
więc $\mathbb EJ_v=q$.
Dla $u\ne v$ izolowanie
obu wierzchołków
wymaga braku $2n-3$
różnych krawędzi
(gdyż krawędź $uv$
byłaby zliczona podwójnie).
Stąd
\[
\operatorname{Cov}(J_u,J_v)
=(1-p)^{2n-3}-(1-p)^{2n-2}
=p(1-p)^{2n-3}.
\]
Sumujemy wartości oczekiwane,
wariancje poszczególnych
funkcji charakterystycznych i kowariancje
uporządkowanych par
różnych wierzchołków.
\end{proof}

\begin{twierdzenie}[Próg spójności]
\label{thm:tg-random-connectivity}
Niech $a_n>0$ dla dostatecznie
dużych $n$, przy czym
$a_n\to\infty$ oraz
$a_n=o(\log n)$.
Zdefiniujmy
\[
p_n^-=\frac{\log n-a_n}{n},
\qquad
p_n^+=\frac{\log n+a_n}{n}.
\]
Dla dostatecznie dużych $n$
obie liczby należą do $[0,1]$.
Wówczas
\[
\mathbb P\bigl(G(n,p_n^-)
\text{ jest spójny}\bigr)\to0,
\qquad
\mathbb P\bigl(G(n,p_n^+)
\text{ jest spójny}\bigr)\to1.
\]
Zatem $(\log n)/n$ jest
progiem spójności
w powyższym sensie.
\end{twierdzenie}
\begin{proof}
Najpierw rozważmy
$p=p_n^-$ i liczbę
izolowanych wierzchołków $Y$.
Ponieważ $p=O((\log n)/n)$,
mamy $np^2\to0$, a zatem
\[
\begin{aligned}
\mathbb EY
&=n(1-p)^{n-1}\\
&=\exp\bigl(\log n-(n-1)p+O(np^2)\bigr)\\
&=\exp(a_n+o(1))\longrightarrow\infty.
\end{aligned}
\]
Dla $q=(1-p)^{n-1}$
wzór z
\ref{prop:tg-random-isolates}
daje
\[
\begin{aligned}
\frac{\operatorname{Var}(Y)}{(\mathbb EY)^2}
&=\frac{1-q}{nq}
+\frac{n-1}{n}\frac{p}{1-p}\\
&\leq\frac1{\mathbb EY}
+\frac{p}{1-p}\longrightarrow0.
\end{aligned}
\]
Lemat \ref{lem:tg-random-second-moment}
pokazuje, że
$\mathbb P(Y>0)\to1$.
Dla $n\geq2$ graf
z wierzchołkiem izolowanym
jest niespójny, zatem
prawdopodobieństwo
spójności dąży do zera.

Niech teraz $p=p_n^+$.
Jeżeli graf jest niespójny,
ma składową o liczności
$k$, gdzie
$1\leq k\leq\lfloor n/2\rfloor$.
Dla ustalonego zbioru
$S\subseteq V$, $|S|=k$,
brak krawędzi między
$S$ i $V\setminus S$
ma prawdopodobieństwo
$(1-p)^{k(n-k)}$.
Wobec tego
\[
\mathbb P(G(n,p)\text{ jest niespójny})
\leq\sum_{k=1}^{\lfloor n/2\rfloor}
\binom nk(1-p)^{k(n-k)}.
\]
Stosujemy
$\binom nk\leq(en/k)^k$
oraz $1-p\leq e^{-p}$.
Dla $k\leq n/\log n$
przy dostatecznie dużym $n$
otrzymujemy
\[
p(n-k)\geq
(\log n+a_n)(1-1/\log n)
\geq\log n+a_n/2-1.
\]
Zatem suma składników
o takich $k$ jest
nie większa niż
\[
\sum_{k\geq1}
\left(\frac{e^2e^{-a_n/2}}{k}\right)^k
\leq\sum_{k\geq1}
\left(e^2e^{-a_n/2}\right)^k
\longrightarrow0.
\]
Dla $n/\log n<k\leq n/2$
mamy $n/k<\log n$
oraz $n-k\geq n/2$, więc
\[
\binom nk(1-p)^{k(n-k)}
\leq\left(
e\log n\,e^{-(\log n+a_n)/2}
\right)^k
\leq\left(\frac{e\log n}{\sqrt n}\right)^k.
\]
Ostatnia suma również
dąży do zera, gdyż jej
podstawa zmierza do zera.
Zatem prawdopodobieństwo
niespójności $G(n,p_n^+)$
dąży do zera.
\end{proof}

\begin{uwaga}[Brak izolowanych wierzchołków
nie jest równoważny spójności]
Zdarzenie braku wierzchołków
izolowanych samo w sobie
nie gwarantuje spójności:
suma rozłączna dwóch
trójkątów nie ma
wierzchołków izolowanych,
lecz jest niespójna.
Dlatego w drugim
kierunku twierdzenia
potrzebne było
oszacowanie wszystkich
przekrojów oddzielających
niepuste składowe.
\end{uwaga}

\begin{przyklad}[Trzy wierzchołki]
\label{ex:tg-random-three-connect}
W grafie $G(3,p)$
warunek spójności
jest równoważny
obecności co najmniej
dwóch spośród trzech
możliwych krawędzi.
Zatem
\[
\mathbb P(G(3,p)\text{ jest spójny})
=3p^2(1-p)+p^3=3p^2-2p^3.
\]
Z kolei prawdopodobieństwo
wystąpienia trójkąta
wynosi $p^3$.
Są to różne własności
i mają różne
prawdopodobieństwa.
\end{przyklad}

\section{Pierwszy moment i występowanie większych klik}
\label{sec:tg-random-cliques}

\begin{stwierdzenie}[Oszacowanie dla kliki]
\label{prop:tg-random-cliques}
Niech $r\geq2$
będzie stałą liczbą
całkowitą, a $Z_r$
liczbą $r$-elementowych
zbiorów wierzchołków
wyznaczających $K_r$
w grafie $G(n,p)$.
Wówczas
\[
\mathbb EZ_r=
\binom nr p^{\binom r2},
\qquad
\mathbb P(Z_r>0)
\leq\binom nr p^{\binom r2}.
\]
Jeżeli $(p_n)_{n\geq1}$ jest ciągiem liczb z $[0,1]$
i $p_n=o(n^{-2/(r-1)})$,
to
$\mathbb P(G(n,p_n)
\text{ zawiera }K_r)\to0$.
\end{stwierdzenie}
\begin{proof}
Dla ustalonego
$r$-elementowego
zbioru wierzchołków
potrzeba obecności
$\binom r2$ niezależnie
wybieranych krawędzi;
prawdopodobieństwo
wynosi zatem
$p^{\binom r2}$.
Liniowość wartości
oczekiwanej daje wzór,
a lemat
\ref{lem:tg-random-first-moment}
daje nierówność.
Jeśli
$p_n=o(n^{-2/(r-1)})$,
to
$n^rp_n^{\binom r2}
=\bigl(n^{2/(r-1)}p_n\bigr)^{\binom r2}
\to0$, co dowodzi
ostatniej tezy.
\end{proof}

Zauważmy, że to
oszacowanie jest
\emph{jednokierunkowe}:
samo obliczenie
wartości oczekiwanej
nie wystarcza
do wykazania istnienia
kliki z prawdopodobieństwem
dążącym do jedności.
Dla trójkątów
potrzebowaliśmy
w tym celu
wariancji i
drugiego momentu.

\begin{twierdzenie}[Dwustronny próg pojawiania się ustalonej kliki]
\label{thm:tg-random-clique-threshold}
Niech $r\geq3$ będzie ustaloną liczbą całkowitą,
a $(p_n)_{n\geq1}$ ciągiem liczb z $[0,1]$.
Niech $Z_{r,n}$ oznacza liczbę kopii $K_r$
na $r$-elementowych zbiorach wierzchołków
losowego grafu $G(n,p_n)$.
Wówczas zachodzą implikacje:
\begin{enumerate}[label=(\roman*)]
\item jeśli $n^{2/(r-1)}p_n\to0$, to
$\mathbb P(Z_{r,n}>0)\to0$;
\item jeśli $n^{2/(r-1)}p_n\to\infty$, to
$\mathbb P(Z_{r,n}>0)\to1$.
\end{enumerate}
Zatem dla każdej ustalonej kliki $K_r$
funkcja $n^{-2/(r-1)}$ wyznacza próg
pojawiania się tej kliki.
\end{twierdzenie}
\begin{proof}
Oznaczmy $b=\binom r2$ oraz
$N_n=\binom nr$. Zgodnie z
\ref{prop:tg-random-cliques},
$\mu_n=\mathbb EZ_{r,n}=N_np_n^b$.
W pierwszym przypadku,
ponieważ $N_n\leq n^r/r!$ i
$r=2b/(r-1)$, mamy
\[
\mu_n\leq\frac1{r!}
\bigl(n^{2/(r-1)}p_n\bigr)^b
\longrightarrow0.
\]
Nierówność Markowa dowodzi (i).

W drugim przypadku dla dostatecznie
dużych $n$ zachodzi $p_n>0$ oraz $n\geq2r$.
Niech $J_S$ oznacza funkcję charakterystyczną zdarzenia,
w którym $r$-elementowy zbiór $S$
wyznacza klikę $K_r$, tak że
$Z_{r,n}=\sum_{S\in\binom Vr}J_S$.
Dla różnych $S,T$ o
$|S\cap T|=j$ mamy
\[
\operatorname{Cov}(J_S,J_T)=
\begin{cases}
0,&j=0,1,\\
p_n^{2b-\binom j2}-p_n^{2b},&2\leq j\leq r-1.
\end{cases}
\]
Istotnie, przy $j\leq1$
obie kliki nie mają wspólnej krawędzi,
więc odpowiednie funkcje charakterystyczne są niezależne.
Przy $j\geq2$ wspólnych krawędzi
jest dokładnie $\binom j2$.
Dla każdego $S$ liczba zbiorów
$T$ o danym $j$ wynosi
$\binom rj\binom{n-r}{r-j}$.
Sumując wariancje i kowariancje
po \emph{uporządkowanych} parach
różnych zbiorów, otrzymujemy
\[
\frac{\operatorname{Var}(Z_{r,n})}{\mu_n^2}
\leq\frac1{\mu_n}
+\sum_{j=2}^{r-1}
\frac{\binom rj\binom{n-r}{r-j}}{\binom nr}
p_n^{-\binom j2}.
\]
Przy ustalonym $r$ i $j$
iloraz współczynników
dwumianowych jest $O(n^{-j})$.
Połóżmy
$\lambda_n=n^{2/(r-1)}p_n\to\infty$.
Dla $2\leq j\leq r-1$ zachodzi
\[
n^jp_n^{\binom j2}
=\lambda_n^{\binom j2}
n^{j(r-j)/(r-1)}\longrightarrow\infty.
\]
Co więcej,
$\mu_n\asymp n^rp_n^b
=\lambda_n^b\to\infty$.
Wszystkie składniki oszacowania
wariancji podzielonej przez $\mu_n^2$
dążą zatem do zera, a liczba
składników sumy jest stała.
Stąd
$\operatorname{Var}(Z_{r,n})/\mu_n^2\to0$.
Na mocy lematu
\ref{lem:tg-random-second-moment}
otrzymujemy
$\mathbb P(Z_{r,n}=0)\to0$,
co dowodzi (ii).
\end{proof}

\begin{przyklad}[Dlaczego próg zależy od wielkości kliki]
\label{ex:tg-random-clique-threshold}
Dla trójkąta $K_3$ próg to $n^{-1}$,
jak w Twierdzeniu
\ref{thm:tg-random-triangle-threshold}.
Dla $K_4$ wynosi on $n^{-2/3}$,
a dla $K_5$ jest równy $n^{-1/2}$.
Na przykład przy
$p_n=n^{-3/4}$ mamy
$n^{2/3}p_n=n^{-1/12}\to0$,
więc w $G(n,p_n)$ z prawdopodobieństwem
dążącym do jedności nie ma $K_4$.
Jednocześnie $np_n=n^{1/4}\to\infty$,
więc trójkąt występuje
z prawdopodobieństwem dążącym
do jedności. Nie ma tu sprzeczności:
występowanie mniejszej kliki
nie pociąga za sobą
występowania większej.
\end{przyklad}

\section{Próg występowania ustalonego podgrafu}
\label{sec:tg-random-fixed-subgraph}

Dla dowolnego grafu zabronionego sama średnia
liczba jego kopii nie zawsze wystarcza do
wyznaczenia progu: przeszkodą może być gęstszy
podgraf. Poniższy wynik uogólnia próg występowania
kliki.

\begin{definicja}[Maksymalna gęstość podgrafu]
\label{def:tg-random-max-density}
Niech $H$ będzie skończonym grafem prostym
nieskierowanym z co najmniej jedną krawędzią.
Definiujemy
\[
\rho(H)=\max_{\substack{F\subseteq H\\ |V(F)|>0}}
\frac{|E(F)|}{|V(F)|},
\]
gdzie maksimum bierzemy po podgrafach, niekoniecznie
indukowanych. Mamy $\rho(H)>0$.
\end{definicja}

\begin{twierdzenie}[Próg ustalonego podgrafu]
\label{thm:tg-random-fixed-subgraph}
Niech $H$ będzie ustalonym skończonym grafem prostym
nieskierowanym z co najmniej jedną krawędzią,
a $p_n\in[0,1]$. Jeśli $n^{1/\rho(H)}p_n\to0$,
to prawdopodobieństwo zawierania $H$ jako podgrafu
w $G(n,p_n)$ dąży do zera. Jeśli
$n^{1/\rho(H)}p_n\to\infty$, to dąży do jedności.
\end{twierdzenie}
\begin{proof}
Oznaczmy $\rho=\rho(H)$. Wybierzmy podgraf $F$
osiągający maksimum, o $v_F$ wierzchołkach
i $e_F>0$ krawędziach. Każda kopia $H$ zawiera
kopię $F$. Liczba iniekcji zbioru $V(F)$ do
$\{1,\ldots,n\}$ jest nie większa niż $n^{v_F}$,
a prawdopodobieństwo, że dana iniekcja zachowuje
wszystkie krawędzie $F$, wynosi $p_n^{e_F}$.
Zatem metodą pierwszego momentu
\[
\mathbb P(H\subseteq G(n,p_n))
\leq n^{v_F}p_n^{e_F}
=(n^{1/\rho}p_n)^{e_F}\to0.
\]

Przyjmijmy teraz $\lambda_n=n^{1/\rho}p_n\to\infty$.
Niech $v=|V(H)|$, $e=|E(H)|$, a $X$ oznacza liczbę
iniekcji $\varphi:V(H)\to\{1,\ldots,n\}$ zachowujących
wszystkie krawędzie $H$. Nie wymagamy zachowania
braku krawędzi. Dla $n\geq v$ mamy
$X=\sum_\varphi J_\varphi$, gdzie $J_\varphi$
jest funkcją charakterystyczną zdarzenia zachowania
krawędzi. Stąd $\mu=\mathbb EX=(n)_v p_n^e$,
gdzie $(n)_v=n(n-1)\cdots(n-v+1)$.
Dla dużych $n$ mamy $(n)_v\geq(n/2)^v$ i
\[
\mu\geq2^{-v}n^{v-e/\rho}\lambda_n^e\to\infty,
\]
ponieważ $e/v\leq\rho$.

Dla dwóch iniekcji $\varphi,\psi$ niech $j$
oznacza liczbę wspólnych wierzchołków ich obrazów,
a $h$ liczbę wspólnych wymaganych krawędzi.
Jeśli $h=0$, funkcje charakterystyczne są niezależne.
Jeśli $h\geq1$, to
$\mathbb E(J_\varphi J_\psi)=p_n^{2e-h}$,
więc kowariancja nie przekracza tej liczby.
Wspólne wierzchołki i wspólne wymagane krawędzie
wyznaczają podgraf obrazu $H$; po przeniesieniu go
przez $\varphi^{-1}$ otrzymujemy podgraf $H$
o $j$ wierzchołkach i $h$ krawędziach.
Zatem $h/j\leq\rho$.

Dla każdego $j$ liczba uporządkowanych par iniekcji
o dokładnie $j$ wspólnych wierzchołkach jest
$O_H(n^{2v-j})$: najpierw wybieramy obraz pierwszej
iniekcji na co najwyżej $n^v$ sposobów, następnie
wybieramy i przyporządkowujemy wspólne pozycje na
liczbę sposobów zależną tylko od $H$, a pozostałe
$v-j$ obrazy wybieramy na co najwyżej $n^{v-j}$
sposobów. Sumujemy kowariancje także dla
$\varphi=\psi$ (wtedy $j=v$, $h=e$).
Z dolnego oszacowania $(n)_v$ wynika
\[
\frac{\operatorname{Var}(X)}{\mu^2}
\leq C_H
\sum_{\substack{1\leq j\leq v,\ 1\leq h\leq e\\h/j\leq\rho}}
\frac{1}{n^jp_n^h},
\]
dla pewnej stałej $C_H$ zależnej tylko od $H$.
Każdy mianownik jest równy
$n^{j-h/\rho}\lambda_n^h\to\infty$,
ponieważ $j-h/\rho\geq0$ i $h\geq1$.
Suma ma skończenie wiele składników, zatem
iloraz wariancji przez $\mu^2$ dąży do zera.
Lemat \ref{lem:tg-random-second-moment}
daje $\mathbb P(X>0)\to1$, co kończy dowód.
\end{proof}

\begin{przyklad}[Gęstszy podgraf wyznacza próg]
Niech $H$ będzie sumą rozłączną $K_4$ i dwóch
wierzchołków izolowanych. Cały graf ma gęstość
$6/6=1$, ale $\rho(H)=6/4=3/2$:
podgraf $K_4$ osiąga tę gęstość, a żaden jego
podgraf nie ma większej. Dodatkowe izolowane
wierzchołki mogą tylko zmniejszyć gęstość.
Próg wynosi więc $n^{-2/3}$, nie $n^{-1}$.
Przy $p_n=n^{-3/4}$ średnia liczba iniekcji
zachowujących krawędzie $H$ jest rzędu
$n^6p_n^6=n^{3/2}\to\infty$, lecz z
prawdopodobieństwem dążącym do jedności
nie ma ani jednej kopii $K_4$, a zatem również $H$.
\end{przyklad}

\section{Probabilistyczne oszacowanie liczb Ramseya}
\label{sec:tg-random-ramsey}

Losowy graf można
również wykorzystać
do konstrukcji
\emph{nie wprost}:
wystarczy wykazać,
że prawdopodobieństwo
istnienia grafu
bez zadanych
konfiguracji jest
dodatnie. Otrzymamy
w ten sposób
dolną granicę
liczby Ramseya
$R(k,k)$ z
rozdziału
\ref{ch:tg-extremal-ramsey}.

\begin{twierdzenie}[Dolne ograniczenie Ramseya
metodą pierwszego momentu]
\label{thm:tg-random-ramsey-bound}
Niech $k\geq2$
oraz $n\geq1$
będą liczbami
całkowitymi.
Jeżeli
\[
\binom nk\,2^{1-\binom k2}<1,
\]
to istnieje graf
prosty na $n$
wierzchołkach,
który nie zawiera
ani kliki $K_k$,
ani $k$-elementowego
zbioru niezależnego.
W szczególności
$R(k,k)>n$.
\end{twierdzenie}
\begin{proof}
Kolorujmy niezależnie
każdą krawędź
$K_n$ na czerwono
lub niebiesko,
każdym kolorem
z prawdopodobieństwem
$1/2$. Dla ustalonego
$k$-elementowego
zbioru $S$ wszystkie
$\binom k2$ krawędzie
muszą mieć jednakowy
kolor, aby otrzymać
jednobarwną $K_k$.
To prawdopodobieństwo
wynosi
$2(1/2)^{\binom k2}
=2^{1-\binom k2}$.
Jeśli $X$ oznacza
liczbę jednobarwnych
klik $K_k$, to
z liniowości wartości
oczekiwanej
\[
\mathbb EX=
\binom nk\,2^{1-\binom k2}.
\]
Na mocy
\ref{lem:tg-random-first-moment}
przy $\mathbb EX<1$
istnieje kolorowanie
z $X=0$. Uznajmy
czerwone krawędzie
za krawędzie grafu $G$,
a niebieskie za jego
dopełnienie. Wówczas
$G$ nie ma $K_k$
ani $k$-elementowego
zbioru niezależnego.
Definicja $R(k,k)$
daje $R(k,k)>n$.
\end{proof}

\begin{wniosek}[Wykładniczy wzrost liczb Ramseya]
\label{cor:tg-random-ramsey-exponential}
Dla każdego
całkowitego $k\geq3$
zachodzi
\[
R(k,k)>
\left\lfloor2^{(k-1)/2}\right\rfloor.
\]
\end{wniosek}
\begin{proof}
Przyjmijmy
$n=\lfloor2^{(k-1)/2}\rfloor$.
Wówczas
\[
\binom nk\,2^{1-\binom k2}
\leq\frac{n^k}{k!}
\,2^{1-\binom k2}
\leq\frac2{k!}<1
\]
dla $k\geq3$.
Stosujemy twierdzenie
\ref{thm:tg-random-ramsey-bound}.
\end{proof}

\begin{uwaga}[Istnienie a konstrukcja]
Dowód daje dodatnie
prawdopodobieństwo
istnienia grafu
o żądanej własności,
ale sam nie wskazuje,
który konkretny
graf ją posiada.
Jest to charakterystyczny
element metody
probabilistycznej.
Do konstrukcji jawnej
potrzebne są dodatkowe
pomysły lub algorytmy.
\end{uwaga}

\section{Podsumowanie i dalsze kierunki}
\label{sec:tg-random-summary}

Porównanie modeli $G(n,p)$ i $G(n,m)$ pokazało,
że warunkowanie względem liczby krawędzi zmienia
niezależność ich wyboru. Sprzężenie monotoniczne
pozwala porównywać prawdopodobieństwa własności
rosnących. Otrzymaliśmy również trzy rodzaje wyników:
\begin{itemize}[leftmargin=2em,itemsep=3pt]
\item wzorów dokładnych
na wartości oczekiwane
i wariancje
liczb krawędzi,
trójkątów oraz
wierzchołków izolowanych;
\item progów
dla wystąpienia
trójkąta
($p\asymp n^{-1}$)
oraz ustalonego podgrafu $H$
($p\asymp n^{-1/\rho(H)}$), a także spójności
($p\asymp(\log n)/n$);
\item dowodu
istnienia
grafów o małych
klikach i małych
zbiorach niezależnych,
prowadzącego do
wykładniczych
dolnych ograniczeń
liczb Ramseya.
\end{itemize}
